\section{Background}

\subsection{Softmax Attention}

\[
\text{Attn}(Q, K, V) = \text{softmax}\!\left(\frac{Q K^\top}{\sqrt{d\_h}}\right) V
\]

with $Q, K, V \in \mathbb{R}^{B \times H \times L \times d_h}$.

\subsection{Linear Attention}

Linear attention replaces softmax with a kernel $\phi$:

\[
\text{Attn}(Q, K, V)\_i = \frac{\sum\_j \phi(q\_i)^\top \phi(k\_j) v\_j}{\sum\_j \phi(q\_i)^\top \phi(k\_j)}
\]

The \textbf{associative} trick:

\[
\text{Attn}(Q, K, V)\_i = \frac{\phi(q\_i)^\top \left[ \sum\_j \phi(k\_j) v\_j^\top \right]}{\phi(q\_i)^\top \left[ \sum\_j \phi(k\_j) \right]}
\]

Computing the bracketed terms is $\mathcal{O}(L \cdot 2F \cdot d_h)$, independent of $L^2$.

\subsection{Random Fourier Features}

\textbf{Bochner's theorem} (Rahimi \& Recht, 2007) states that a shift-invariant 
kernel $k(x, y) = k(x - y)$ can be represented as the Fourier transform of 
a probability distribution $p(\omega)$:

\[
k(x - y) = \mathbb{E}\_\omega [e^{i \omega^\top (x - y)}]
\]

Approximating this expectation with $F$ samples:

\[
k(x - y) \approx \frac{1}{F} \sum\_{i=1}^{F} e^{i \omega\_i^\top (x - y)}
\]

Using Euler's formula, the real part is:

\[
k(x - y) \approx \frac{1}{F} \sum\_{i=1}^{F} \cos(\omega\_i^\top (x - y))
= \phi(x)^\top \phi(y)
\]

where $\phi(x) = \frac{1}{\sqrt{F}} [\cos(\omega_i^\top x), \sin(\omega_i^\top x)]_{i=1}^F$.

\subsection{The Yat-Kernel}

The \textbf{Yat-Kernel} (Zhang et al., 2024) modifies the RFF to:

\[
\phi(x) = \frac{1}{F^{1/4}} [\cos(x \Omega) \,\|\, \sin(x \Omega)]
\]

where $\Omega \in \mathbb{R}^{d_h \times F}$ is a random projection matrix.

\textbf{Properties:}

\begin{itemize}
\item $2F$ output dimension.
\item $\phi(x)^\top \phi(y) = \frac{1}{\sqrt{F}} [\cos(x\Omega)^\top \cos(y\Omega) + \sin(x\Omega)^\top \sin(y\Omega)]$.
\item Using $\cos(a)\cos(b) + \sin(a)\sin(b) = \cos(a - b)$, this equals
  $\frac{1}{\sqrt{F}} \sum_i \cos(\omega_i^\top (x - y))$.

\end{itemize}
The Yat-Kernel is thus a scaled cosine-similarity kernel, which is 
\textbf{positive} (needed for stable linear attention).

\subsection{Spherical Normalization}

Given Q, K, we project them onto the unit sphere:

\[
\hat{Q} = \frac{Q}{\|Q\|\_2 + \epsilon}, \quad
\hat{K} = \frac{K}{\|K\|\_2 + \epsilon}
\]

\textbf{Properties:}

\begin{itemize}
\item $\|\hat{Q}\|_2 = 1$ for all tokens and heads.
\item Attention scores depend only on \textbf{direction}, not magnitude.
\item Stabilizes the linear kernel: $\phi(\hat{K})$ has bounded magnitude.

\end{itemize}
\subsection{Adaptive Gating}

A learnable gate blends two streams:

\[
g = \sigma(W\_g X + b\_g)
\]

\[
O = g \odot O^\text{dense} + (1 - g) \odot O^\text{linear}
\]

\textbf{Initialization.} $W_g = 0$, $b_g = 5.0$ gives $g \approx 
0.993$ at step 0.

\subsection{LoRA Recovery}

\[
\Delta W = \frac{\alpha}{r\_L} B A
\]

with $B = 0$ initially.

\subsection{Notation}

\begin{table}[H]
\centering
\small
\begin{tabular}{ll}
\toprule
\textbf{Symbol} & \textbf{Meaning} \\
\midrule
$Q, K, V$ & Query, Key, Value \\
$\hat{Q}, \hat{K}$ & Spherical-normalized Q, K \\
$\phi$ & Yat-Kernel feature map \\
$\Omega$ & RFF projection \\
$F$ & Number of RFF features \\
$g$ & Adaptive gate \\
$L, d_h$ & Sequence length, head dim \\
\bottomrule
\end{tabular}
\end{table}

